\chapter{BỐI CẢNH}

\section{Đặt vấn đề} % Sẽ tự động là 1.1 nếu không có cấu trúc khác, hoặc bạn dùng cấu trúc dưới đây:
Đột quỵ vẫn là một trong những nguyên nhân hàng đầu gây tử vong và tàn tật trên toàn thế giới. Tỷ lệ đột quỵ gia tăng đòi hỏi các chiến lược phòng ngừa hiệu quả và chẩn đoán kịp thời. Việc sử dụng các công nghệ tiên tiến như học máy [1--3] có thể cải thiện đáng kể dự đoán nguy cơ và khả năng nhận diện sớm những bệnh nhân cần được chăm sóc. Hiểu được cơ chế của các rối loạn mạch máu não là chìa khóa để phát triển các chiến lược phòng ngừa và điều trị đột quỵ hiệu quả, qua đó giảm thiểu hậu quả của bệnh và cải thiện chất lượng cuộc sống của bệnh nhân.
\\
Việc sử dụng trí tuệ nhân tạo trong bối cảnh này mở ra những triển vọng mới để cải thiện thực hành y khoa. Tích hợp các mô hình học máy như XGBoost hoặc mạng nơ-ron cho phép dự đoán chính xác nguy cơ đột quỵ dựa trên các đặc trưng riêng của từng bệnh nhân [4,5]. Điều này bao gồm cả các yếu tố nguy cơ truyền thống lẫn các tham số mới được xác định thông qua phân tích lượng dữ liệu lớn. Nhờ đó, các chuyên gia y tế có thể đánh giá nguy cơ và triển khai các biện pháp phòng ngừa kịp thời hiệu quả hơn. Các thuật toán học máy như XGBoost thể hiện mức độ chính xác cao khi giải quyết những bài toán phức tạp, bao gồm cả dự báo y khoa. Việc tối ưu các thuật toán này cho những mục tiêu cụ thể, trong đó có dự đoán nguy cơ đột quỵ, là một bước quan trọng để nâng cao hiệu quả của các quyết định y tế.

\section{Vai trò của xử lý dữ liệu và giảm chiều (PCA)}
Giảm tác động của các bất thường và nhiễu trong dữ liệu là một bước then chốt trong xử lý dữ liệu, bởi chất lượng dữ liệu ảnh hưởng trực tiếp đến độ chính xác và độ tin cậy của mô hình học máy. Các bất thường như ngoại lệ hoặc dữ liệu sai có thể làm sai lệch đáng kể kết quả phân tích, dẫn đến các kết luận không chính xác hoặc thậm chí các khuyến nghị có hại [6]. Phân tích thành phần chính (PCA) giúp xác định và loại bỏ các bất thường này bằng cách tập trung vào những đặc trưng chính và giàu thông tin nhất của tập dữ liệu [7,8]. Điều này không chỉ cải thiện chất lượng dữ liệu mà còn nâng cao độ tin cậy của kết quả, bởi các mô hình được huấn luyện trên dữ liệu đã làm sạch có thể đưa ra dự báo chính xác và ổn định hơn.

\noindent Ngoài ra, PCA làm giảm độ phức tạp của mô hình, điều đặc biệt quan trọng trong các ứng dụng y khoa nơi các quyết định có thể ảnh hưởng đến sức khỏe bệnh nhân [9]. Bằng cách giảm số lượng biến cần phân tích, PCA đơn giản hóa quá trình huấn luyện mô hình và dẫn đến xử lý dữ liệu nhanh hơn. Sự đơn giản của các mô hình sau khi áp dụng PCA khiến chúng dễ hiểu và dễ diễn giải hơn đối với các chuyên gia y tế vốn có thể không có kiến thức chuyên sâu về thống kê hay học máy. Điều này quan trọng vì bác sĩ và nhân viên lâm sàng cần hiểu cách thức và lý do một dự báo cụ thể được đưa ra trước khi sử dụng trong thực hành. Hơn nữa, giảm số lượng biến giúp tránh quá tải thông tin, từ đó nâng cao khả năng hiểu các yếu tố nguy cơ nền tảng. Vì vậy, ứng dụng PCA trong nghiên cứu và thực hành y khoa là một bước quan trọng hướng tới tối ưu hóa quá trình ra quyết định và tăng hiệu quả của can thiệp y tế.

\noindent Tuy nhiên, thuật toán tuần tự ở giai đoạn tiền xử lý có thể tiêu tốn nhiều thời gian do độ phức tạp tính toán cao, bởi PCA truyền thống đòi hỏi tài nguyên đáng kể để tính các trị riêng và vector riêng của ma trận hiệp phương sai, đặc biệt nghiêm trọng đối với các tập dữ liệu lớn. Khi quá trình xử lý chỉ diễn ra trên một lõi CPU, tài nguyên tính toán sẵn có không được tận dụng, dẫn đến độ trễ. Vì vậy, trong nghiên cứu này, nhóm tác giả đề xuất sử dụng các công nghệ tính toán song song hiện đại, cụ thể là OpenMP, cho phép áp dụng phương pháp đề xuất trên bất kỳ hệ thống máy tính đa lõi nào [10–12]. Điều này sẽ tăng hiệu quả xử lý dữ liệu nhờ tối ưu hóa việc sử dụng tài nguyên tính toán, từ đó giảm thời gian thực thi thuật toán và nâng cao hiệu năng.

\section{Nhu cầu về trí tuệ nhân tạo giải thích được (XAI)}
Việc áp dụng PCA có thể vẫn là một ``hộp đen'', gây khó khăn cho việc diễn giải và hiểu những yếu tố nào đã đóng góp vào quá trình giảm chiều dữ liệu. Điều này có thể làm giảm niềm tin của các chuyên gia y tế và nhà nghiên cứu đối với mô hình, vì họ có thể không hiểu tại sao một số đặc trưng lại tác động đến kết quả mạnh hơn các đặc trưng khác. Vì vậy, cần tăng cường bước này bằng các kỹ thuật XAI [13], giúp xác định những đặc trưng có ảnh hưởng lớn nhất đến kết quả PCA, nhờ đó các nhà nghiên cứu và chuyên gia y tế hiểu rõ hơn cách các đặc tính cụ thể của dữ liệu hình thành kết quả phân tích. Tương tự nghiên cứu [14] đề xuất một cách tiếp cận để tăng cường bảo mật điện toán đám mây bằng cách phân tích hành vi người dùng và đánh giá độ tin cậy của họ, nhóm tác giả sử dụng các phương pháp trí tuệ nhân tạo giải thích được nhằm cải thiện khả năng diễn giải và tính minh bạch của các mô hình dự đoán nguy cơ đột quỵ.

\section{Tổng quan các phương pháp học máy hiện có cho dự đoán đột quỵ}
Dự đoán đột quỵ là một khía cạnh quan trọng của y học, bởi chẩn đoán kịp thời có thể làm giảm đáng kể hậu quả của bệnh và cải thiện chất lượng cuộc sống của bệnh nhân. Nhiều phương pháp học máy khác nhau được sử dụng trong nghiên cứu khoa học, mỗi phương pháp có ưu điểm và nhược điểm riêng. Bảng~\ref{tab:methods} trình bày so sánh các phương pháp học máy hiện có cho dự đoán đột quỵ.

\begin{longtable}{p{2.3cm}p{4.5cm}p{3.7cm}p{4.0cm}}
\caption{So sánh các phương pháp học máy hiện có cho dự đoán đột quỵ}\label{tab:methods}\\
\toprule
\textbf{Phương pháp} & \textbf{Mô tả} & \textbf{Ưu điểm} & \textbf{Nhược điểm}\\
\midrule
\endfirsthead
\toprule
\textbf{Phương pháp} & \textbf{Mô tả} & \textbf{Ưu điểm} & \textbf{Nhược điểm}\\
\midrule
\endhead
Hồi quy Logistic & Mô hình hóa xác suất đột quỵ dựa trên các yếu tố nguy cơ. Sử dụng hàm sigmoid và gradient descent để xây dựng mô hình. Đánh giá hiệu năng mô hình khi có và không có regularization. & Độ chính xác cao (trên 95\%), có thể cải thiện thông qua regularization và dễ triển khai. & Độ chính xác hạn chế khi quan hệ phi tuyến; nhạy với lựa chọn tham số.\\
\hline %
Cây quyết định & Một tập các cây quyết định được huấn luyện trên các tập con dữ liệu khác nhau, sử dụng biểu quyết để đưa ra quyết định cuối cùng. & Trực quan hóa quyết định, dễ diễn giải. & Có xu hướng overfitting nếu không regularization.\\
\hline %
Random Forest & Tổng hợp kết quả từ nhiều cây quyết định, làm giảm nguy cơ overfitting. & Độ chính xác cao (96\%), độ tin cậy tốt. & Có thể chậm trên các tập dữ liệu lớn.\\
\hline %
Naive Bayes & Phân loại dựa trên giả định các đặc trưng độc lập. & Đơn giản, hiệu quả trong nhiều bài toán phân loại. & Độ chính xác đạt 82\%; hạn chế khi tồn tại quan hệ phức tạp giữa các đặc trưng.\\
\hline %
k-NN & Phân loại các quan sát mới dựa trên các láng giềng gần nhất trong tập huấn luyện. & Đơn giản và rõ ràng. & Khó mở rộng, nhạy với lựa chọn tham số $k$.\\
\hline %
SVM & Sử dụng các hàm kernel để xử lý phân bố phi tuyến. & Hiệu quả cao trên dữ liệu nhiều chiều. & Nhạy với lựa chọn tham số, gặp khó khăn với tập dữ liệu lớn.\\
\hline %
Học sâu (CNN) & Sử dụng mạng nơ-ron tích chập (CNN) để phân tích ảnh y khoa. & Độ chính xác cao trong phát hiện các mẫu phức tạp. & Đòi hỏi tập dữ liệu lớn và tài nguyên tính toán mạnh.\\
\hline %
Mạng nơ-ron nhân tạo (ANN) & Mô hình học máy sử dụng resampling, tránh rò rỉ dữ liệu, lựa chọn đặc trưng và các kỹ thuật diễn giải như permutation importance và LIME để dự đoán đột quỵ. & Khả năng diễn giải cao với LIME, resampling và lựa chọn đặc trưng hiệu quả; độ chính xác dự báo cao (95\%). & Phụ thuộc vào kiểm định trên tập dữ liệu bên ngoài và cần tối ưu liên tục để đạt hiệu năng tốt hơn.\\
\hline %
XGBoost & Thuật toán tăng cường gradient kết hợp nhiều phương pháp để đạt kết quả tốt hơn. & Hiệu quả dự báo và khả năng diễn giải cao; độ chính xác trên 97\%. & Khó cấu hình tham số, yêu cầu tài nguyên tính toán.\\
\bottomrule
\end{longtable}

\section{Các nghiên cứu liên quan}
Công trình [15] so sánh các thuật toán như hồi quy logistic, phân loại cây quyết định, random forest và voting classifier; random forest đạt độ chính xác cao nhất, khoảng 96\%. Nghiên cứu [16] đánh giá hiệu quả các thuật toán phân loại, đặc biệt Naive Bayes, đạt độ chính xác khoảng 82\%, nhưng dữ liệu bị giới hạn ở một khu vực địa lý duy nhất, hạn chế khả năng khái quát hóa.

\noindent Nghiên cứu [17] và [18] tập trung vào stacking và lựa chọn đặc trưng tự động; [17] cho thấy phương pháp stacking đạt AUC 98.9\%, trong khi [18] nhấn mạnh tầm quan trọng của dữ liệu chất lượng cao. Các nghiên cứu [19] và [20] sử dụng XGBoost, đạt độ chính xác 97.56\% và AUC 0.8595 tương ứng, cho thấy tiềm năng áp dụng lâm sàng ở cấp cá nhân.

\noindent Công trình [21] sử dụng XGBoost để dự đoán sự xuất hiện của bệnh, đạt độ chính xác cao (97\%) khi dự đoán trường hợp không đột quỵ, nhưng độ chính xác dự đoán đột quỵ chỉ khoảng 20\%; nghiên cứu này cũng gặp dữ liệu thiếu ở biến "bmi". Công trình [22] cho thấy Random Forest đạt độ chính xác cao nhất ($\sim 96\%$), nhưng bị giới hạn bởi kích thước mẫu. Công trình [23] nghiên cứu CNN và LSTM, cho thấy các phương pháp học máy truyền thống trong một số trường hợp vẫn vượt học sâu.

\noindent Các tác giả của [24] đề xuất dự đoán đột quỵ tim bằng ANN với độ chính xác 95\%, áp dụng permutation importance và LIME, nhưng không sử dụng giảm chiều dữ liệu (PCA) hay tính toán song song — đây là hai điểm mà nghiên cứu hiện tại bổ sung so với [24]. Công trình gần đây nhất [25] trình bày tối ưu XGBoost bằng ensemble và tối ưu Bayes, cho hiệu năng cao nhưng chưa nhấn mạnh việc tích hợp cơ chế giải thích để tăng độ tin cậy trong ứng dụng lâm sàng.

\noindent Tổng quan tài liệu cho thấy những thách thức chính vẫn là khả năng diễn giải thấp của mô hình học máy và độ phức tạp khi xử lý các tập dữ liệu y tế lớn có nhiều đặc trưng. Phần lớn phương pháp hiện đại dùng các thuật toán mạnh để cải thiện độ chính xác nhưng chưa quan tâm đầy đủ đến giảm chiều dữ liệu và giải thích đóng góp của từng đặc trưng vào kết quả cuối cùng [26,27].

\section{Mục tiêu và đóng góp của nghiên cứu}
Nghiên cứu này nhằm cải thiện quá trình dự đoán nguy cơ đột quỵ bằng cách tích hợp thuật toán XGBoost với phương pháp PCA tối ưu và triển khai các phương pháp XAI nhằm tăng tính minh bạch và khả năng diễn giải của kết quả phân tích.
Đóng góp chính của phương pháp đề xuất gồm:

\begin{itemize}
  \item Kết hợp XGBoost với PCA ở giai đoạn tiền xử lý để giảm chiều và cải thiện cấu trúc dữ liệu, qua đó tăng độ chính xác của mô hình.
  \item Song song hóa PCA bằng công nghệ OpenMP, giảm đáng kể thời gian xử lý các tập dữ liệu lớn — quan trọng trong các ứng dụng y tế nơi tốc độ ra quyết định có thể mang tính sống còn.
  \item Tích hợp XAI vào quá trình PCA, làm tăng tính minh bạch và khả năng hiểu mô hình, hỗ trợ phát hiện bất thường và nhiễu trong dữ liệu.
\end{itemize}

\noindent Do đó, tính mới của phương pháp nằm ở việc áp dụng tích hợp các phương pháp học máy, công nghệ tính toán và công cụ giải thích, cho phép tạo ra các mô hình dự báo nguy cơ đột quỵ chính xác hơn, nhanh hơn và dễ diễn giải hơn, có ý nghĩa lớn đối với sự phát triển của ngành chăm sóc sức khỏe.
